\documentclass[11pt]{article}
\usepackage[T1]{fontenc}
\usepackage{lmodern,amsmath,amssymb,amsthm,mathtools}
\usepackage[margin=1in]{geometry}
\usepackage{booktabs,longtable,array,microtype,xcolor}
\usepackage{xurl}
\usepackage[colorlinks=true,linkcolor=blue!45!black,urlcolor=blue!55!black,citecolor=blue!45!black]{hyperref}
\hypersetup{pdftitle={Pell Congruences, Shared Positivity, and Smaller Diophantine Consistency Equations},pdfauthor={AI-assisted research prepared for Vladimir Reshetnikov},pdfsubject={Explicit Diophantine equations equivalent to syntactic inconsistency of ZFC},pdfkeywords={Diophantine equations, ZFC, consistency, Pell equations, CRT}}
\usepackage{enumitem,fancyhdr}
\setlist{itemsep=3pt,topsep=5pt}
\setlength{\emergencystretch}{2em}
\newtheorem{theorem}{Theorem}[section]
\newtheorem{lemma}[theorem]{Lemma}
\newtheorem{proposition}[theorem]{Proposition}
\newtheorem{corollary}[theorem]{Corollary}
\theoremstyle{remark}\newtheorem{remark}[theorem]{Remark}
\newcommand{\Z}{\mathbb Z}
\newcommand{\N}{\mathbb N}
\newcommand{\Pos}{\operatorname{Pos}}
\newcommand{\Con}{\operatorname{Con}}
\newcommand{\Refute}{\operatorname{Ref}}
\newcommand{\SOS}{\operatorname{SOS}}
\newcommand{\pin}{\texttt{cf4a5218a5a4e4358d3f5314151b6b05d4aa384e}}
\pagestyle{fancy}\fancyhf{}\lhead{Pell congruences and Diophantine consistency}\rhead{\thepage}
\setlength{\headheight}{14pt}
\title{Pell Congruences, Shared Positivity, and\ Smaller Diophantine Consistency Equations}
\author{Research contribution prepared for Vladimir Reshetnikov\\\small AI-assisted derivation and independent internal audits}
\date{23 October 2026}
\begin{document}
\maketitle
\begin{abstract}
We give explicit fixed-coefficient polynomials whose absence of integer
solutions is equivalent, by ordinary mathematics formalizable in ZFC, to
syntactic consistency of ZFC. Relative to the pinned wave-10 ProveIt
snapshot, the principal new integer tradeoffs include a quartic in 37
unknowns, an octic in 26, degree 10 in 24, degree 16 in 22, degree 20 in
20, degree 32 in 18, and degree 88 in 17. The main new compression replaces two auxiliary
quotients by one coordinate using a congruence which the preliminary Pell
rank argument already proves. This avoids a circular appeal to the exact
rank which those quotients are intended to certify. A parity-forced
change of the first Pell root, an exact transport projection, and compatibility between existing
determinant and residue guards yield further reductions. We also give
eleven-positive-variable degree 269 and twelve-positive-variable degree
78 equations. A separately certified nine-positive-variable companion
trades more arithmetic gates for lower degree. All coefficients, literal
polynomials, builders and finite certificates accompany the article.
These are improvements to an inspected collection, not a proof of global
minimality or a new machine-independent calibration. The native compiler,
specified machine/refutation equivalence and named Pell rank theorems
remain explicit dependencies; no new end-to-end proof-assistant
certificate is claimed.
\end{abstract}
\tableofcontents

\section{Results, comparison and logical meaning}
\label{sec:results}

The comparison is frozen at ProveIt commit \pin. This is the README
reconciliation after wave 10, including packages 68--75. The repository's
own review lists 38 unrestricted integer variables at degree four, 32 at
degree six, 27 at degree eight, and 26 at degree sixteen. It also retains
an older 24-integer degree-1936 form. These are domain-qualified results:
the older sextic in 25 \emph{positive} variables and octic in 20 positive
variables do not describe unrestricted-integer zero sets
\cite{repo,review}.

The following new forms all use exactly the stated domain, actual fixed
integer coefficients, and one final polynomial equation.
\begin{center}
\begin{tabular}{rrrrl}
\toprule
Domain & Unknowns & Exact degree & Residuals & Main mechanism\\
\midrule
$\Z$ & 37 & 4 & 24 & Determinant plus residue orientation\\
$\Z$ & 31 & 6 & 18 & Monic projection\\
$\Z$ & 26 & 8 & 13 & Monic projection\\
$\Z$ & 24 & 10 & 11 & Centered root and transport projection\\
$\Z$ & 23 & 12 & 11 & Shared positivity\\
$\Z$ & 22 & 16 & 10 & Shared positivity and transport projection\\
$\Z$ & 20 & 20 & 8 & Retain the rank as one coordinate\\
$\Z$ & 19 & 28 & 7 & CRT and transport projection\\
$\Z$ & 18 & 32 & 6 & Further product projection\\
$\Z$ & 17 & 88 & 5 & Eliminate the remainder coordinate\\
\bottomrule
\end{tabular}
\end{center}
Here ``residuals'' means that the single polynomial is the sum of their
squares. The six-residual, eighteen-variable form is displayed in full
in Section~\ref{sec:explicit}. Named polynomial registers in that display
are abbreviations, not additional quantified variables.

The integer results are accompanied by the following positive-domain
tradeoffs. An operation is one live binary addition, subtraction or
multiplication; squares and multiplication by constants are charged.
Fixed numerals and reuse of a previously computed register are free.
This is an arithmetic-circuit convention, not bit complexity.
\begin{center}
\begin{tabular}{lrrr}
\toprule
Construction & Positive unknowns & Exact degree & Operations\\
\midrule
Wave-10 source 70 & 11 & 325 & 113\\
New saturated transport & 11 & 269 & 111\\
Wave-10 source 69 & 12 & 82 & 102\\
New adaptive saturated merger & 12 & 78 & 102\\
Wave-10 source 71 & 9 & 1,379,730 & 478\\
New full-radical majorant & 9 & 1,158,228 & 726\\
\bottomrule
\end{tabular}
\end{center}
The nine-variable companion improves degree but increases evaluator size.
It does not replace the separate 443-operation local nine-variable target
reported by the active research. The eleven-variable result improves
both degree and the stated operation count.

\subsection{The exact metamathematical assertion}
Fix the recursively specified ZFC calculus used by the repository's
machine calibration and let $\Refute_{\mathrm{ZFC}}(p)$ mean that $p$ codes
a finite derivation of a contradiction. Put
\[
 \Con(\mathrm{ZFC})\;\equiv\;\neg\exists p\in\N\;\Refute_{\mathrm{ZFC}}(p).
\]
The construction establishes, with the explicitly named inherited
interfaces, the arithmetic theorem
\begin{equation}
 \mathrm{ZFC}\vdash
 \left(\exists\boldsymbol x\in\Z^n\;P(\boldsymbol x)=0\right)
 \longleftrightarrow
 \exists p\in\N\;\Refute_{\mathrm{ZFC}}(p).       \label{eq:logical}
\end{equation}
Negating both sides gives the requested no-solution equivalence. Positive
variables replace $\Z^n$ in the positive companions. No coefficient is
defined by deciding consistency, waiting for a refutation, or consulting
an oracle. No assertion that ZFC is consistent is needed in proving the
biconditional.

Ordinary mathematical formalizability and a supplied formal proof are
different deliverables. The new arguments below use integer arithmetic,
finite recurrences, classical number theory and induction, all available
in ZFC. This package does not contain a newly kernel-checked Lean/Rocq
proof or a serialized derivation of \eqref{eq:logical} in the
repository's object-language calculus. Its finite computer checks
certify coefficients and algebra, not the unbounded semantic interfaces.

\subsection{What is new and what is reused}
The determinant guard is package 71's construction; residue orientation
is package 70's; the nearest-square joint positivity lemma is package
68's. We prove their simultaneous compatibility, then add a lossless
CRT parametrization at the correct stage of the signed Pell argument.
The elementary CRT identity is not asserted to be a new number-theoretic
theorem. Its use to remove these two source quotients, with a proof that
the needed congruence is available before auxiliary rank recovery, is
the new structural reduction. Replacing the slack coordinate by the
native $C$ coordinate also makes the transport row monic in $F$, enabling
one more exact projection. Centering and further projection produce
the new literal equations and their degree certificates. No external
priority claim beyond the inspected sources is made.

The main integer family uses neither Jones's universal polynomial nor
Sun's ten-integer reduction. It still uses the calibrated native encoding
of K215. A separate analysis of direct proof-checker and finitely
axiomatized class-theory routes appears in Section~\ref{sec:alternatives};
those alternatives are not counted as completed new equations here.

\section{Fixed coefficients and the inherited source}
\label{sec:source}

The integer forms use the complete source-34 tuple
\[
 (\beta,K_0,M_C,M_F),\qquad \beta=2^{1064921}-1,
\]
with $M_F$ already shifted by $\beta$. Two additional fixed scales are
\begin{equation}
 A_F=2^{1113},\qquad A_C=2^{666104}.           \label{eq:scales}
\end{equation}
These six numbers are constants of the polynomial, not witness variables.
The package stores all six actual integers as gzip-compressed minimal
unsigned little-endian byte strings, with hashes of the decompressed
integers. Their total uncompressed size is 615,883 bytes. This is a real
coefficient-size cost, not hidden in the variable count.

\subsection{A finite recipe, not a parameter family}
The retained \texttt{coefficients/source34/build\_coefficients.py} defines
the four native constants by a finite recipe from the literal 430
transitions of K215. It derives and verifies the ordered 943-tile table.
With $\rho=2^{121}$, block length 8801, and $B=\rho^{8801}$, the recipe
constructs finite selector/copy weights and a finite convolution
polynomial $D_C$. It sets
\begin{align*}
 \beta&=B-1,& K_0&=D_C+B\rho,\\
 M_C&=\text{the specified low mask},&
 M_F&=\text{the specified high mask}+B-1.
\end{align*}
The masks, tile order, support positions and every finite sum are literal
in that short builder and its retained coefficient proof. The supplied
payloads are checked against the builder, including the mask shift.
Consequently the displayed formulas specify one closed polynomial once
these attached fixed numerals are substituted. The builder never runs
the blank machine to completion. The compact source-44 tuple used later
for positive companions is kept separately and is never mixed into
these integer coefficients. The nine-positive-variable companion uses
the complete source-31 tuple instead: its transport shift is
$D_R=\rho^{42}$, whereas source 34 uses $D_R=\rho$. Equality of their
$\beta$ values does not make the other coefficients interchangeable.

\subsection{The source theorem we import}
Packages 64 and 70 supply the signed native/Pell source used here
\cite{source64,source70}. The following is the exact logical boundary.
The source has arithmetic identities which, after eliminating its
monic registers, read
\begin{align*}
 q&=\beta J+1,& C&=q-F-\alpha,& X&=qW,\\
 Y&=s_0q^3,& E&=XY,& a&=(X+1)Y,\\
 H&=4a+3,&\Delta&=(a+2)^2-1,& k&=\eta+\zeta,\\
 c&=kY+\eta,& P&=k+E(c-\eta),&D&=X+ac+H\gamma,\\
 I&=i_0c^2,& I^2&=\Delta(f^2-1),&
 V&=fo-c=jc-R.
\end{align*}
The remaining abbreviations are
\[
 R=((q-1)(q-F)+\alpha)(q^2-1)+(M_C+qM_F)J,
 \qquad k=R+1+hE.
\]
Here $s_0,i_0$ are odd, $W$ is a nonzero multiple of four, and the
source enforces the common-sign requirements
\[
 F\alpha>0,\quad FC>0,\quad\eta\zeta>0.
\]
Its four remaining norm/transport equations are
\begin{align}
 \tau_0^2-P^2+Pk-1&=0,                         \label{eq:oldfirst}\\
 D^2-\Delta c^2-1&=0,                         \label{eq:oldmain}\\
 I^2V^2-(I^2-1)y_0^2-1&=0,                   \label{eq:oldaux}\\
 (K_0+W)C-F-z(q-1)&=0.                        \label{eq:oldtransport}
\end{align}
Package 70 removes the separate $X=2u_0^2$ witness by imposing
$J=3J_3+1$ and $W=4(3w_3+1)$. Its signed power theorem first proves
$X=\pm2^R$ with odd $R$; then $X\equiv2\pmod3$ forces the positive
sign. Completing the source gives a finite refutation, and any finite
refutation gives canonical source witnesses.

One part of this proof must be distinguished from its conclusion. Before
using either congruence for $V$ or the auxiliary norm
\eqref{eq:oldaux}, the first/main norms, the index relation and the
ordinary strong equation yield
\begin{equation}
 |c|=\psi_b(p),\quad |f|=\chi_b(m),\quad
 p\mid m,\quad c\ne0,\quad f\ne0,
 \qquad b=|a+2|\ge2.                         \label{eq:earlyrank}
\end{equation}
The proof appears in package 70, Sections 4.1--4.2, and the retained
ordinary strong-rank proof. The exact rank $p=R$ is obtained only later.
This distinction is essential in Section~\ref{sec:crt}.

\section{Compatible sign reductions and a 37-variable quartic}
\label{sec:sign}

For integers define
\[
 \Pos(x,y,z)=1+x^2+y^2+z(z+1).
\]
It is always positive and represents every positive integer. For the
latter assertion apply the three-square theorem to $4n-3$; one of the
three squares is odd and the other two are even. The familiar
representation of nonnegative integers by $x^2+y^2+z^2+z$ is also
formally documented in \cite{threeSquares}.

\begin{lemma}[Determinant sign recovery]
For fixed $A_F,A_C\ge2$, put
\[
 U=A_FF,\quad W_C=A_CC,\quad q=F+\alpha+C,\quad
 G=UW_C-(q-U-W_C)^2.
\]
If $G>0$, then $F,C,\alpha$ have one common nonzero sign.
\end{lemma}
\begin{proof}
The inequality implies $UW_C>0$ and
$|q-U-W_C|<\sqrt{UW_C}\le(|U|+|W_C|)/2$.
Thus $q$ has their common sign and
$|q|>(|U|+|W_C|)/2\ge|F|+|C|$. Subtracting $F+C$ gives the same
strict sign for $\alpha$. No sign of $q$ has been assumed.
\end{proof}

\begin{proposition}[Simultaneous canonical completeness]
The determinant guard with \eqref{eq:scales} and the two affine residue
restrictions of package 70 are simultaneously satisfiable on the
canonical source witnesses of every finite refutation.
\end{proposition}
\begin{proof}
The source-34 canonical family has $q=2^t$, $t=1064921N$ odd,
$R\equiv3\pmod4$ and $R>t$. Its original coordinates satisfy
\[
 J=\frac{2^t-1}{2^{1064921}-1},\qquad
 w_{\rm old}=2^{R-t-2}.
\]
Therefore $J\equiv1\pmod3$ and $w_{\rm old}\equiv1\pmod3$, so
$J=3J_3+1$ and $w_{\rm old}=3w_3+1$ have integer inverse coordinates.

Package 71 proves, for this same rank-tuned family,
\[
 \frac q2\le A_FF<\frac{4q}{5},\qquad
 \frac q2\le A_CC<\frac{3q}{5}.               \label{eq:canonicalbox}
\]
Its proof is retained in the archive. The highest occupied $C$ digit is
at exponent 3296 in radix $2^{121}$; all lower support is at most 2195.
The leading $F$ digit at exponent 8791 has coefficient
105534389824824194647382298625. The rank controller changes a dummy
at exponent 1063, below these leading positions, and the carry bounds
cover every permitted dummy change. Thus the estimates are compatible
with the parity/rank tuning, rather than results for two unrelated
witness families.

Write $U=q/2+x$, $W_C=q/2+y$. Then
$0\le x<3q/10$, $0\le y<q/10$, and
\[
 G=(q/2)^2+(q/2)(x+y)-x^2-xy-y^2>3q^2/25>0.
\]
Choose a ternary $\Pos$ representation of this positive integer.
\end{proof}

Consequently package 70's two guards on $F\alpha$ and $FC$, with six
private coordinates, can be replaced by $G=\Pos$ with three. Soundness
restores the two former guards \emph{before} invoking the signed power
theorem. Completeness uses the one compatible family just proved. The
result is a quartic in 37 integers, the sum of 24 quadratic residuals.
Its coefficient of $\tau_0^4$ is one.

The direct-coordinate alternative replaces only $X$ by $3X_3-1$ and
also has 37 unknowns. It is included as
\texttt{quartic37\_direct.json}; it has a shorter expanded expression.
The affine-$J,w$ version is used for the monic projections below,
because it retains a monic definition of $X$.

\section{The new CRT compression}
\label{sec:crt}

\begin{lemma}[One parameter for two source congruences]
\label{lem:crt}
Let $c,f$ be nonzero integers with $c\mid f^2-1$. For every integer $R$,
the simultaneous congruences
\[
 V\equiv-c\pmod f,\qquad V\equiv-R\pmod c
\]
hold exactly for the integers
\begin{equation}
 V=fcT-f^2R-c,\qquad T\in\Z.                 \label{eq:crt}
\end{equation}
For fixed $c,f,R$, this parametrization is one-to-one.
\end{lemma}
\begin{proof}
Since $f^2\equiv1\pmod c$, we have $\gcd(f,c)=1$. Write
$V=fo-c$. The second congruence implies
\[
 c\mid f(o+fR)=V+c+f^2R,
\]
so $c\mid o+fR$. Set $T=(o+fR)/c$.
Conversely, from \eqref{eq:crt} take
\begin{equation}
 o=cT-fR,\qquad
 j=fT-1-\frac{(f^2-1)R}{c}.                  \label{eq:restoreoj}
\end{equation}
The second expression is an integer by hypothesis, and these give
$V=fo-c=jc-R$. The formulas are inverse. Signs do not enter the proof.
\end{proof}

\begin{lemma}[The CRT hypothesis is available early]
Every tuple satisfying the preliminary signed source equations used in
\eqref{eq:earlyrank} satisfies $f^2\equiv1\pmod{c^2}$.
\end{lemma}
\begin{proof}
The Pell identity is
$\chi_b(m)^2-1=(b^2-1)\psi_b(m)^2$.
Since $p\mid m$, divisibility of the Pell sequence gives
$\psi_b(p)\mid\psi_b(m)$. Use \eqref{eq:earlyrank} and square the
divisibility. This does not use the auxiliary norm, either quotient
$o,j$, or the later conclusion $p=R$.
\end{proof}

\begin{theorem}[Lossless quotient elimination]
In the residue-oriented, determinant-guarded source, the two rows
$V=fo-c$ and $V=jc-R$, and their private coordinates $o,j$, may be
replaced by the one row \eqref{eq:crt} and one coordinate $T$ without
changing existence of an integer solution.
\end{theorem}
\begin{proof}
For soundness first recover the signs and preliminary rank facts from
the equations which do not use $o,j,V,y_0$. The preceding lemma supplies
the hypothesis of Lemma~\ref{lem:crt}. Restore $o,j$ by
\eqref{eq:restoreoj}; the rest of the source, including its auxiliary
norm, is now literally available. Its final rank/calibration proof
applies. For completeness take any source solution and use the forward
CRT formula for $T$. All other coordinates stay unchanged.
\end{proof}

There is no stronger requirement $f\equiv\pm1\pmod c$. Nor do we
impose $f^2\mid V+c$: that would retain only a subprogression and need a
new completeness argument. The coefficient of $cT$ in
\eqref{eq:crt} is $f$, and this distinction is what makes the reduction
lossless. The only displayed division is in the \emph{proof} of the
inverse map, after its integrality is proved; the final polynomial has
no division.

\section{Centered roots and shared positivity}
\label{sec:centering}

Two algebraic identities lower the degrees of the projections. First,
put $P=k+E(c-\eta)$ and $\tau_0=\tau+E(c-\eta)$. Then
\begin{equation}
 \tau_0^2-P^2+Pk-1
 =\tau^2+(2\tau-k)E(c-\eta)-1.               \label{eq:centerfirst}
\end{equation}
This integral triangular change is inherited from the recent source
projections. Second, put $y_0=V+u$ in the auxiliary norm:
\begin{equation}
 I^2V^2-(I^2-1)y_0^2-1
 =V^2+(1-I^2)u(2V+u)-1.                      \label{eq:centeraux}
\end{equation}
When $V$ is itself a polynomial, this eliminates its unnecessarily
repeated leading square from the product with $I^2$.

\begin{lemma}[A parity-forced half-centered root]
\label{lem:half}
For integers $P,k$, integer roots of \eqref{eq:oldfirst} are in
bijection with integer roots of
\begin{equation}
 t^2+2t(2P-k)+k^2-4=0.                       \label{eq:half}
\end{equation}
The maps are $t=2\tau_0-2P+k$ and
$\tau_0=(t+2P-k)/2$.
\end{lemma}
\begin{proof}
Substitution gives four times the old residual. Conversely the new
residual differs from $(t-k)^2$ by $4tP-4$. Hence a zero has
$(t-k)^2\equiv0\pmod4$ and $t-k$ even. Its inverse is integral and
restores the old norm. No parity promise is an added constraint.
\end{proof}

After $\eta=c-kY$ has been substituted, the usual centered first row
has degree 24. Formula \eqref{eq:half} has degree 18. This is useful in
the intermediate twenty-variable, degree-44 unpaired form, although the
shared guard dominates the degree of the eventual seventeen-variable
form.

\begin{lemma}[Shared strict positivity; package 68]
\label{lem:pair}
For integers $x,y$, put
\[
 \mathcal G(x,y,s)=xy-(s^2-x-y+1)^2.
\]
Then $x,y>0$ if and only if $\mathcal G(x,y,s)>0$ for some integer $s$.
\end{lemma}
\begin{proof}
A positive value implies $xy>0$. If $x=-a,y=-b$ with $a,b\ge1$,
the subtracted square is at least $(a+b+1)^2>ab$, a contradiction.
Conversely set $N=x+y-1\ge1$. Then $xy\ge N$. Choose the closest
integer square $s^2$ to $N$; writing $m=\lfloor\sqrt N\rfloor$ shows
that either $s=m$ or $m+1$ gives $|s^2-N|\le m$, with square error
strictly below $N$ unless $N$ is square, when the error is zero.
Thus $(s^2-N)^2<N\le xy$.
\end{proof}

The two equations $G=\Pos$ and $\eta\zeta=\Pos$, with six private
coordinates, may therefore be replaced by
\begin{equation}
 G\eta\zeta-(s^2-G-\eta\zeta+1)^2
 =1+p_1^2+p_2^2+p_3(p_3+1),                  \label{eq:shared}
\end{equation}
with four private coordinates. This step is an existence equivalence,
not an identity between the old and new sums of squares. Its proof and
credit are retained explicitly \cite{source68}.

\section{A complete polynomial in eighteen integers}
\label{sec:explicit}

\subsection{A transport coordinate which can be eliminated}
Replace the independent coordinate $\alpha$ by $C=q-F-\alpha$.
This is an integral triangular bijection, with inverse
$\alpha=q-F-C$. The transport equation is now monic in $F$:
\begin{equation}
 F=(K_0+W)C-z(q-1).                         \label{eq:transportprojection}
\end{equation}
Substitute this expression everywhere and delete the transport row and
the independent coordinate $F$. The packing index becomes
\begin{equation}
 R=\bigl(q(q-F)-C\bigr)(q^2-1)+(M_C+qM_F)J. \label{eq:newR}
\end{equation}
These are polynomial maps in both directions on the corresponding zero
sets. No divisibility, sign, rank or consistency assumption is needed
for this projection. Its degree cost depends on which arithmetic
coordinates have already been eliminated: here $R$ has degree five and
$G$ degree four, while the leading forms of $k$ and the shared guard
remain unchanged.

\subsection{The polynomial}
This section specifies the principal compact form without hidden
quantifiers. Supply exactly the following eighteen integers, with no
positivity restrictions:
\begin{equation}
 j,C,z,f,h,r,v,w,\tau,\eta,u,\gamma,T,c,s,p_1,p_2,p_3.
 \label{eq:variables19}
\end{equation}
Use the six fixed coefficients from Section~\ref{sec:source}. Define
polynomial abbreviations in the following order:
\begin{gather}
 J=3j+1,\qquad q=\beta J+1,\qquad W=4(3w+1),\nonumber\\
 F=(K_0+W)C-z(q-1),\quad X=qW,\quad Y=(2v+1)q^3,\nonumber\\
 E=XY,\qquad a=E+Y,\qquad H=4a+3,\qquad\Delta=a^2+H,\nonumber\\
 R=\bigl(q(q-F)-C\bigr)(q^2-1)+(M_C+qM_F)J,\nonumber\\
 k=R+1+hE,\qquad\zeta=k-\eta,\qquad I=(2r+1)c^2,\nonumber\\
 V=fcT-f^2R-c,\qquad D=X+ac+H\gamma,\nonumber\\
 G=A_FF\,A_CC-(q-A_FF-A_CC)^2,\qquad Z_\eta=\eta\zeta.
                                                        \label{eq:regs19}
\end{gather}
The six residuals are
\begin{align}
 r_c&=c-kY-\eta,\nonumber\\
 r_0&=\tau^2+(2\tau-k)E(c-\eta)-1,\nonumber\\
 r_m&=D^2-\Delta c^2-1,\nonumber\\
 r_s&=I^2-\Delta(f^2-1),\nonumber\\
 r_a&=V^2+(1-I^2)u(2V+u)-1,\nonumber\\
 r_g&=GZ_\eta-(s^2-G-Z_\eta+1)^2
                 -1-p_1^2-p_2^2-p_3(p_3+1).             \label{eq:rows19}
\end{align}
The requested single Diophantine equation is
\begin{equation}
 \boxed{P_{18}=r_c^2+r_0^2+r_m^2+r_s^2+r_a^2+r_g^2=0.}
                                                        \label{eq:p19}
\end{equation}
All exponents are fixed. There are no variable powers, rational
coefficients, implicit sign conditions or additional quantified registers.
The JSON witness names \texttt{J3}, \texttt{C\_native},
\texttt{v\_scale}, \texttt{w3}, \texttt{y\_gap} and
\texttt{pair\_root} correspond to $j,C,v,w,u,s$ here.

\begin{theorem}
\label{thm:p19}
The polynomial \eqref{eq:p19} has exact total degree 32. With the
fixed coefficients above, it has an integer zero if and only if the
specified ZFC calculus has a finite refutation. In particular ZFC proves
that it has no integer zero if and only if $\Con(\mathrm{ZFC})$.
\end{theorem}
\begin{proof}[Soundness]
A zero makes all six residuals vanish. Lemma~\ref{lem:pair} applied to
$r_g$ gives $G>0$ and $\eta\zeta>0$. The determinant lemma restores
the two earlier common-sign guards. Restore the polynomial registers in
\eqref{eq:regs19}, $\alpha=q-F-C$, the old first root
$\tau_0=\tau+E(c-\eta)$, and $y_0=V+u$.
Then $r_c$ supplies the original $c=kY+\eta$, and every preliminary
signed source equation is present. Its early strong-rank conclusion
\eqref{eq:earlyrank} gives $c,f\ne0$ and $c\mid f^2-1$ without using
either missing quotient. Lemma~\ref{lem:crt} now restores the integer
quotients $o,j_{\rm aux}$ in \eqref{eq:restoreoj}. Equations
\eqref{eq:centerfirst}--\eqref{eq:centeraux} recover all original norms,
and \eqref{eq:transportprojection} supplies the original transport
equation identically. Thus the guarded,
residue-oriented source is fully restored. Its signed power and native
calibration theorems yield a finite refutation.
\end{proof}
\begin{proof}[Completeness]
From a finite refutation choose the common canonical source witnesses of
Section~\ref{sec:sign}, including both affine residue restrictions and
$G>0$. Compress the two source quotients by Lemma~\ref{lem:crt}.
Project all monically defined arithmetic coordinates, use
\eqref{eq:transportprojection}, and choose the
integral centered roots. Since $G$ and $\eta\zeta$ are positive,
Lemma~\ref{lem:pair} and the ternary $\Pos$ representation supply
$s,p_1,p_2,p_3$. Every row of \eqref{eq:rows19} vanishes. All choices
are finite arithmetic constructions from a hypothetical refutation;
none assumes such a refutation actually exists.
\end{proof}
\begin{proof}[Exact degree]
The leading forms in the input degree are
\[
 q_1=3\beta j,\quad X_2=36\beta jw,\quad
 Y_4=54\beta^3j^3v,\quad E_6=1944\beta^4j^4vw.
\]
Thus $[k]_7=hE_6$ and $[Z_\eta]_8=\eta hE_6$.
The guard has degree 16, with leading form $-\eta^2h^2E_6^2$.
The other five residuals have degrees $11,14,14,14,14$.
Consequently the degree-32 leading form of $P_{18}$ is precisely
\[
 \eta^4h^4E_6^4
 =1944^4\beta^{16}j^{16}v^4w^4h^4\eta^4,
\]
whose coefficient is positive for the actual $\beta$. This is an
exact-degree proof after coefficient specialization.
\end{proof}

\subsection{A seventeen-variable equation of degree 88}
Delete $\eta$ from \eqref{eq:variables19}, define $\eta=c-kY$, and
delete $r_c$. Put $P=k+E(c-\eta)$ and replace $r_0$ by
\[
 r_0^{\rm half}=\tau^2+2\tau(2P-k)+k^2-4.
\]
Keep the other four rows. The sum of squares of these five rows is an explicit equation
$P_{17}=0$ in seventeen unrestricted integers. Lemma~\ref{lem:half}
supplies the integral inverse root, so the proof of
Theorem~\ref{thm:p19} applies.

The leading forms are now
\[
 [\eta]_{11}=-hE_6Y_4,\quad
 [\zeta]_{11}=hE_6Y_4,\quad
 [Z_\eta]_{22}=-h^2E_6^2Y_4^2.
\]
The shared guard has degree 44 with leading form
$-h^4E_6^4Y_4^4$. Its square has degree 88 and positive leading
coefficient. All other rows have degree at most 18. Hence the claimed
degree is exact. The gain in variables comes with a higher degree;
neither $P_{17}$ nor $P_{18}$ dominates the other in both measures.

\subsection{Twenty variables at degree twenty}
There is a particularly useful alternative between the lower-degree
forms and $P_{18}$. Add two independent integer coordinates $E$ and
$R_*$ to the eighteen coordinates in \eqref{eq:variables19}.
Let $R_{\rm native}$ denote the polynomial on the right of
\eqref{eq:newR}. Add the two residuals
\[
 r_R=R_*-R_{\rm native},\qquad r_E=E-XY.
\]
In every other abbreviation use $R=R_*$, keep $E$ independent, and set
$a=E+Y$ as before. Take the sum of squares of these two rows and the six
rows in \eqref{eq:rows19}. This defines $P_{20}$ in twenty integers.
The two new rows uniquely restore $E$ and $R$, so the zero sets are
polynomially isomorphic after forgetting these two coordinates. The
consistency equivalence follows immediately from Theorem~\ref{thm:p19}.

The residual degrees are now $5,6,6,4,10,10,10,8$, in the order
$r_R,r_E,r_c,r_0,r_m,r_s,r_a,r_g$. In particular $k$ has degree two and
$V$ degree three. The degree-ten part of the auxiliary row is
\[
 [r_a]_{10}=-8r^2c^5ufT+8r^2c^4uf^2R_*.
\]
No other row involves both $r$ and $u$. Thus the coefficient of
$r^4c^8u^2f^4R_*^2$ in $P_{20}$ is exactly $64$, proving exact degree
20 independently of the large fixed coefficients. This illustrates why
choosing \emph{which} polynomial coordinate to retain matters: retaining
the native index $R$ is better here than retaining the auxiliary $V$.

\section{The rest of the integer tradeoff family}
\label{sec:family}

The exact packets include intermediate forms so that the reductions can
be inspected locally. From the 37-variable quartic, the six monic
eliminations
\[
 P=k+E(c-\eta),\quad N=q(L+1),\quad X=qW,\quad
 C_2=c^2,\quad F_2=f^2,\quad g=y_0^2-V^2
\]
together with \eqref{eq:centerfirst} give 31 variables and cubic
residuals. Eliminating eleven lift coordinates in total from the
37-variable source, including
$L=q^2-1$, $N=q^3$, $B_{\rm aux}=I^2$, $D$, $\Delta$, and
$\zeta=1-V+cj+hE-\eta$, gives 26 variables and quartic residuals.
The defining rows vanish identically under the supplied restoration
maps. Their inverse simply restores each deleted coordinate.

In that 26-variable system, set $V=fo-c$ and $y_0=V+u$.
The auxiliary row becomes \eqref{eq:centeraux}, of degree five;
every other remaining row has degree at most four. There are 25
variables. The two sign arguments have degrees at most two and three,
so shared positivity has degree six and removes two private coordinates.
This gives the 23-integer, degree-12 form.

For the CRT family it is advantageous to restore
$\zeta=R+1+hE-\eta$ from the packing-index row instead of the
auxiliary quotient row. Retain $V$ initially, impose
$V-fcT+f^2R+c=0$, and substitute $I=(2r+1)c^2$.
Before sharing the sign guards there are 24 variables and eleven rows,
with maximum residual degree eight. Then proceed as follows.
\begin{center}
\begin{tabular}{lrrr}
\toprule
Unpaired source change & Variables & Residuals & SOS degree\\
\midrule
Retain $V,E,a$ & 24 & 11 & 16\\
Eliminate $V=fcT-f^2R-c$ & 23 & 10 & 26\\
Eliminate $a=E+(2v+1)q^3$ & 22 & 9 & 26\\
Eliminate $E=XY$ & 21 & 8 & 28\\
Eliminate $\eta=c-kY$ and half-center & 20 & 7 & 44\\
\bottomrule
\end{tabular}
\end{center}
Sharing the two remaining sign guards changes these to 22/20, 21/26,
20/26, 19/32, and 18/88 in variables/degree. The 21/26 form is dominated
by 20/26 and is not a primary release target. The other listed target
packets are included, except that the degree-88 endpoint is emitted
directly in its smaller transport-projected form.

Finally apply the transport projection
\eqref{eq:transportprojection}--\eqref{eq:newR}. It sends 25/10 to
24/10, 23/12 to 22/16, 20/26 to 19/28, 19/32 to 18/32, and 18/88 to
17/88. Each deletes exactly one coordinate and one row. The 23/12 form
remains a useful tradeoff because projection increases its degree.
The 20/26 form is retained as a derivation step and is dominated by the
rank-register construction $P_{20}$ above. The compact packets in
\texttt{transport/} retain ordered
polynomial abbreviations rather than expanding a large guard needlessly;
these are ordinary fixed-degree polynomials, with no additional unknowns.

When $a$ is retained, the literal builder uses the lower-degree rows
$E(X+1)-Xa$ and $c-\eta-k(a-E)$, together with $a-E-Y$.
After $a$ is eliminated they become $E-XY$ and $c-\eta-kY$.
These are exact row operations because
\begin{align*}
 E(X+1)-Xa&=(E-XY)-X(a-E-Y),\\
 c-\eta-k(a-E)&=(c-\eta-kY)-k(a-E-Y).
\end{align*}
There is no unsupported substitution of a rational quotient.

The highest-degree certificates are also transparent. Before deleting
$V$, the centered auxiliary row has leading term
$-4r^2c^4u^2$ among its degree-eight terms. After deleting $V$, its
degree-thirteen pure-$j$ term is
$648\beta^4r^2c^4uf^2j^4$. Its square has coefficient
$419904\beta^8>0$. This persists when $a$ is eliminated, proving exact
degree 26 for the twenty-variable shared-guard endpoint. For every
expanded packet the independent finite audit additionally certifies a concrete
nonzero specialized leading coefficient modulo 1,000,000,007. For the
transport packets it certifies the leading coefficient of a specified
univariate specialization, with every actual fixed coefficient reduced
modulo the same prime. A second algorithm verifies this coefficient
using finite differences and the proved degree upper bound.
Such a nonzero modular residue is an exact nonvanishing certificate,
not a randomized numerical estimate.

\input{../positive/positive_reductions.tex}

\section{A majorant for the complete Pell radical}
\label{sec:full-radical}

A separate construction improves the lowest inspected degree with nine
strictly positive unknowns. Its degree is
\[
 \boxed{1\,158\,228},
\]
compared with package~71's $1\,379\,730$, a reduction of $16.0540\%$
\cite{source71}.
Its delivered evaluator uses $726$ binary operations, compared with $478$;
this is a degree--size tradeoff. All four coefficient payloads are the
coherent original31 tuple. The new polynomial does not add a witness or
change the native compiler.

\subsection{The complete radical and the denominator that can be cleared}

For $A\ge3$, $C,i\ge1$, and $t\ge3$, put
\begin{align*}
 D&=(A^2-1)C^2+1,& e&=i(4A-5)C,\\
 F&=(D-1)e^2D^2+1,& L&=CD+1-A,\\
 G&=A+FL,& H&=C+Ft,\\
 I&=(G^2-1)H^2+1.
\end{align*}
All these registers are positive, and $G>1$. Regard $A$ as a formal
indeterminate over $\mathbb Q(C,i,t)$. The radical $\sqrt{DFI}$ has a unique
formal Laurent expansion with leading term
$2^{10}C^{24}i^5tA^{23}$. Its polynomial part provides the construction
\begin{equation}
 T_0=2^{16}C^7[A^{\ge0}]\sqrt{DFI}.
 \label{eq:full-root-definition}
\end{equation}
The expression $T_0$ is an integer polynomial with $654$ monomials. Its
complete table is the \texttt{majorant} field of
\path{nine_full_radical/laurent_candidate.json}; the coordinate order is
$(A,C,i,t)$. This table is an ordinary division-free definition.
Formula~\eqref{eq:full-root-definition} explains how it was discovered.

The exact polynomial
\begin{equation}
 \Delta_0=T_0^2-2^{32}C^{14}DFI
 \label{eq:full-root-defect}
\end{equation}
has $A$-degree $22$. Together with the leading term of $T_0$, this also
proves the formal-polynomial-part assertion: dividing the difference of
the squares by their formal sum leaves only negative powers of $A$.
The factor $C^7$ clears variable denominators without raising the
$A$-degree.

\paragraph{Complete positivity certificate.}
Under
\[
 A=x+3,\qquad C=y+1,\qquad i=z+1,\qquad t=u+3,
\]
$T_0$ has exactly $9\,216$ nonzero coefficients and $\Delta_0$ exactly
$32\,637$. Every one is a positive integer, including the constant
coefficients. The archive contains both complete expansions and an
independent verifier that reconstructs the radical and checks every
coefficient. Hence, throughout the stated real domain,
\begin{equation}
 T_0>2^{16}C^7\sqrt{DFI}>0.
 \label{eq:full-root-bound}
\end{equation}
This is a finite algebraic proof. No tests on a bounded selection of
integer inputs replace the coefficient identities.

The improvement comes from approximating the whole radical. The preceding
seventh-order construction approximated $\sqrt{DF}\,G$ and multiplied the
error by the high-degree factor $H$. Taking the polynomial part of the
complete radical prevents that multiplication of the approximation error.

\subsection{The final polynomial and its correctness}

Retain the source and Pell registers of package~71, with positive unknowns
\[
 b,g,h,s,w,C,i,j,n.
\]
The source has convolution height $1873$ and maximum coordinate weight
$116$. Its full literal circuit, all kernels, and the fixed coefficient
generator and compressed integer payloads are included in the companion
directory. In particular,
\begin{gather*}
 U=N^2w,\quad Y=N^2s,\quad M=RY,\quad P=2M^2U,\\
 K=R+1+h(P-1),\quad A=M(U+1),\quad B_1=2R+1,\quad W_1=bw,\\
 J=4A-5,\quad Q=3W_1C-2(W_1^2-1),\quad t=B_1+(j-1)C,\\
 a_0=(P^2-1)K^2+1,\quad V_0=PK,\\
 m=K^2(4\operatorname{cap}-g)
      -16\operatorname{cap}(C-KY)^2,\\
 \Omega=Jm,\qquad E_0=J(n+W_1^2)+Q.
\end{gather*}
Define
\begin{gather*}
 S=2^{18}C^7,\qquad T=4T_0(A,C,i,t),\qquad a_1=S^2DFI,\\
 \Theta=E_0-\Omega(V_0+T),\\
 \boxed{\displaystyle
 P_9=(\Theta^2-\Omega^2(a_0+a_1))^2
                -4\Omega^4a_0a_1.}
\end{gather*}
All displayed registers are abbreviations, not additional unknowns.
The authoritative evaluator is
\path{nine_full_radical/laurent9_positive.json}.

The unchanged source estimates give $R\ge1$, $A\ge3$, and hence
$t\ge3$, unconditionally on positive inputs. Thus
$T>\sqrt{a_1}$ before any source zero has been decoded. Also
$V_0^2-a_0=K^2-1>0$. If $m=0$ the norm is $E_0^4>0$; otherwise its real
factorization gives signed roots $u_0^2=a_0$, $u_1^2=a_1$ and
\[
 E_0/J=m(V_0+T+u_0+u_1).
\]
The parenthesis is strictly positive. Therefore $m>0$ before the
square-root integrality argument. The inherited margin argument gives
$C>K$, $C>B_1+W_1$, and $DFI>a_0$. In particular $a_1>a_0$ and the
rational sum $u_0+u_1$ is nonzero. The identity
\[
 2(u_0+u_1)u_0=(u_0+u_1)^2+a_0-a_1
\]
makes both roots rational and then integral. Since $(u_1/S)^2=DFI$,
the rational-root argument also makes $u_1/S$ an integer. The factor
equation recovers $J\mid Q$, and the existing ratio theorem supplies the
encoded positive source zero.

Conversely, choose the inherited ratio witnesses from a source zero,
take the positive integral roots $r_0=\sqrt{a_0}$ and
$r_1=\sqrt{DFI}$, and set
\[
 n=m(V_0+T+r_0+Sr_1)-W_1^2-Q/J.
\]
It is an integer. The inherited bound
$W_1^2+Q/J<2D$, together with $T>S\sqrt{DFI}>SD>2D$ and $m\ge1$,
shows $n>0$. One norm factor now vanishes. Thus the positive existence
statements are equivalent in both directions. Combining with the
specified native compiler and refutation calibration gives the desired
ZFC-provable consistency equivalence. These inherited semantic results
are not being presented as newly formalized proofs.

\subsection{Exact degree and evaluator cost}

For this unchanged source let
\[
 r=\deg R=7963,\quad q=\deg N=2115,
 \quad a=\deg A=16425,\quad \omega=\deg\Omega=82360.
\]
The leading $T_0$ monomial is $2^{26}A^{23}C^{31}i^5t$; the highest
weighted monomial of~\eqref{eq:full-root-defect} is
\[
 141664530202624\,A^{22}C^{42}i^{10}t^2.
\]
Therefore
\[
 \deg T=23a+r+36=385774,\qquad
 \deg(T^2-a_1)=22a+2r+52=377328.
\]
Write $\delta_0=K^2-1$, $\delta_1=T^2-a_1$, and
\[
 B_* =E_0^2-2E_0\Omega(V_0+T)
                        +\Omega^2(\delta_0+\delta_1).
\]
The exact identity
\[
 P_9=B_*^2+4\Omega^2V_0TB_*
                  +4\Omega^4(\delta_0T^2+a_0\delta_1)
\]
has four candidate degrees
\[
 1084096,\qquad1149781,\qquad1158228,\qquad821246.
\]
The unique dominant term is $4\Omega^4\delta_0T^2$.
On the positive diagonal its leading coefficient is
\[
 2^{166}a_b^{578828}c_b^{276}>0,
 \qquad a_b=2^{2134939},\quad c_b=2^{17079501}.
\]
Consequently the asserted degree is attained by the actual fixed
coefficient instance.

The $726$-operation evaluator uses a $110$-term correction to the old
majorant, in the existing registers $J,C,i,t,D$:
\[
 T=64C^7(T_7-640i-4096)H
                   -\operatorname{Corr}(J,C,i,t,D).
\]
The correction table and the exact identity check are included. There
are $430$ multiplications, $265$ additions, and $31$ subtractions;
coefficient leaves and register reuse are free under the same convention
as the comparison source.

The new degree limit is now set by the first-root defect $K^2-1$. A tighter
simultaneous bound for both norm roots is a concrete further research
question, as are cheaper evaluators of the correction and an application
to the ten-integer finalizer. None of those additional improvements is
claimed here.


\section{Verification, interpretation and integration}
\label{sec:verification}

The release is designed to make each mathematical claim inspectable.
The integer builder reads the source packets as inert arithmetic
expressions, replaces the two guard rows, emits the projected and CRT
families, and records exact counts, liveness and degree certificates.
The parser accepts only integer literals, names, additions,
subtractions, products and fixed nonnegative powers. An independent
checker uses a separately written sparse polynomial ring and does not
import the builder. It verifies the source bridge, row operations,
successive projections, parity identity, factored formulas and actual
coefficient specializations. A separate transport checker proves the
local inverse identities, compares every retained row and definition,
and verifies degree by finite differences independently of the builder's
univariate polynomial algorithm.

The coefficients are independently regenerated from the retained finite
machine/tile recipe. The positive and nine-variable directories contain
their own offline verifiers and coherent fixed coefficient payloads.
The accompanying receipts state the checks actually performed. All new
release checks use explicit failures rather than Python assertions;
an optimized-Python run is also supported.

It is neither possible nor useful to validate a consistency equation by
exhaustively searching for a zero. A finite-box search would say little
about \eqref{eq:logical}. The checks here instead establish exact finite
claims: the emitted polynomial is the proved polynomial, its coefficients
are the specified numerals, its transforms are algebraically correct,
and its counts and leading forms are accurate. The proofs of unbounded
Pell rank, native source soundness and machine/refutation calibration
remain the explicitly retained mathematical inputs.

The recommended starting point for integration is the eighteen-variable
equation \eqref{eq:p19}: it has only six residuals, moderate degree 32,
and a direct reconstruction proof. For a lower degree at the cost of
one coordinate, use the nineteen-variable degree-28 form; at the cost
of two, use the twenty-variable degree-20 form. For fewer
integer variables, use the seventeen-variable degree-88 form. For a
quartic target, use the direct 37-variable variant. There is no
mathematically justified total ordering which calls one of these
``the simplest'' without specifying a metric.

\section{Alternative encodings and further research}
\label{sec:alternatives}

The current snapshot contains a direct syntactic proof-checker route and
an executable direct quartic construction. Its native-independent
internal ZFC derivation and useful measured size remain unfinished in
the reviewed sources \cite{alternatives}. We did not treat an abstract
appeal to DPRM, an unexpanded compiler call, or a finite local transition
relation as a smaller completed consistency equation.

A promising separate direction replaces the many axiom-schema branches
of a direct ZF/ZFC checker by a fixed finite first-order axiomatization
of a suitable class theory. A conservative equiconsistency theorem,
exact proof calculus, and its internal arithmetic encoding would still
be required. Counting fewer local constructors alone does not measure
the global polynomial: syntax validation and unbounded proof aggregation
can dominate. The accompanying alternative-route note makes these
remaining requirements explicit.

The following questions arise directly from the proved reductions.
\begin{enumerate}
\item \textbf{Formalize the quotient lemma at its actual dependency point.}
  Prove the preliminary $p\mid m$ interface, then $f^2\equiv1\pmod{c^2}$,
  before importing any exact auxiliary index. This is a smaller and
  cleaner formal target than immediately formalizing the whole polynomial.
\item \textbf{Search for other inverses already supplied by a norm.}
  The CRT reduction uses $f$ as its own inverse modulo $c$. Other norm
  equations may supply polynomial inverses for scale or transport
  congruences and eliminate further quotient coordinates.
\item \textbf{Jointly optimize projection and degree.}
  Eliminating a monic coordinate can substantially increase the degree
  of a shared positivity guard. An optimizer should keep separate
  scores for arithmetic rows and sign rows, with verified inverse maps.
\item \textbf{Reduce the seventeen-variable guard degree.}
  Its degree 88 comes entirely from applying shared positivity after
  eliminating $\eta$. A different representation of that sign condition
  could improve the last tradeoff without changing any machine encoding.
\item \textbf{Lower fixed coefficient height.}
  The determinant proof uses special source-34 leading anchors. A
  genuine proof for the whole smaller compact tuple could shrink the
  fixed numerals; substituting unrelated ports is not justified.
\item \textbf{Improve full-radical evaluation.}
  The new nine-variable degree gain costs 248 additional gates relative
  to source 71. More factorization or a shorter certified majorant
  would be needed for simultaneous improvement in both measures.
\item \textbf{Complete a distinct calibration.}
  A direct proof-checker or finite-theory route should deliver its
  literal coefficients and the internal consistency biconditional,
  together with actual variable and operation counts. It would then be
  meaningfully comparable with the completed native-source equations.
\end{enumerate}

The present contribution gives substantial new low-degree tradeoffs and
an independently reusable quotient-elimination argument. It does not
prove an absolute minimum, remove the inherited native calibration, or
establish external novelty beyond the inspected repository and cited
mathematics.

\begin{thebibliography}{99}
\raggedright
\bibitem{repo} V. Reshetnikov, \emph{ProveIt: Con(ZFC) in Diophantine form},
  snapshot \pin, \href{https://github.com/VladimirReshetnikov/ProveIt/tree/cf4a5218a5a4e4358d3f5314151b6b05d4aa384e/Computability/DiophantineConsistency/con-zfc-diophantine-forms}{repository collection}.
\bibitem{review} ProveIt, \emph{PROPOSAL\_REVIEW\_WAVE10.md}, same snapshot;
  retained comparison and proof-status boundary.
\bibitem{source64} ProveIt package 64, \emph{Structural reductions},
  \texttt{integer41/PROOF41.md}, \texttt{PROOF44.md}, and literal quartic.
\bibitem{source68} ProveIt package 68, \emph{Joint positivity},
  \texttt{integer\_system/PAIR\_GUARD\_PROOF.md}.
\bibitem{source69} ProveIt package 69, \emph{Cones and Pell saturation},
  \texttt{positive/PROOF.md} and \texttt{integer/PROOF.md}.
\bibitem{source70} ProveIt package 70, \emph{Residue orientation and root
  transfer}, \texttt{integer/PROOF.md}, \texttt{eleven/PROOF.md}, and
  literal source packets.
\bibitem{source71} ProveIt package 71, \emph{Calibrated determinant guards},
  \texttt{proofs/DETERMINANT\_GUARD.md} and
  \texttt{nine\_positive/PROOF.md}.
\bibitem{source74} ProveIt package 74, \emph{Centered Pell gaps},
  the adaptive carrier and cofactor argument.
\bibitem{alternatives} ProveIt, \emph{ALTERNATIVE\_ENCODINGS.md}, same snapshot.
\bibitem{threeSquares} A. Danilkin and L. Chevalier,
  \emph{Three Squares Theorem}, Archive of Formal Proofs, 2023,
  \url{https://isa-afp.org/entries/Three_Squares.html}.
\bibitem{matiyasevich} Yu. V. Matiyasevich, \emph{Matiyasevich theorem},
  Scholarpedia 3(7), 7095 (2008),
  \url{https://www.scholarpedia.org/article/Matiyasevich_theorem}.
\end{thebibliography}
\end{document}
